% =====================================================================
% ejercicio09.tex
% Responsabilidad única: comparar, para los tres mejores k según F1 del
% Ejercicio 8, el efecto de cambiar la métrica de distancia.
% =====================================================================
\subsection{Ejercicio 9 --- Distancia Euclidiana vs. Manhattan}

Se comparan los tres mejores valores de $k$ según F1 obtenidos en el
Ejercicio 8: $k=1,3,5$.

\begin{table}[H]
\centering
\begin{tabular}{rlrrrr}
\toprule
$k$ & Distancia & Accuracy & Precision & Recall & F1 \\
\midrule
1 & Euclidiana & 0.9044 & 0.2496 & 0.2144 & \textbf{0.2307} \\
1 & Manhattan  & 0.9054 & 0.2518 & 0.2105 & 0.2293 \\
3 & Euclidiana & 0.9276 & 0.3978 & 0.1621 & 0.2303 \\
3 & Manhattan  & 0.9283 & 0.4000 & 0.1466 & 0.2146 \\
5 & Euclidiana & 0.9329 & 0.4936 & 0.1357 & 0.2128 \\
5 & Manhattan  & 0.9329 & 0.4921 & 0.1237 & 0.1977 \\
\bottomrule
\end{tabular}
\caption{Euclidiana vs. Manhattan para $k \in \{1,3,5\}$.}
\end{table}

\begin{enumerate}

\item \textbf{¿La métrica modifica el desempeño?}
Sí, pero de forma mucho más discreta que la estandarización o que $k$,
y sin un patrón del todo consistente, pues para $k=1$ y $k=3$ Manhattan
da Accuracy y Precision levemente superiores pero Recall y F1
levemente inferiores a Euclidiana, mientras que para $k=5$ ocurre lo
contrario y la Euclidiana iguala o supera a Manhattan en las cuatro
métricas (obsérvese la Precision de $0.4936$ para Euclidiana y no
$0.4921$ como en la Manhattan), aunque en ningún caso la diferencia entre
métricas supera $1.5$ puntos porcentuales, muy por debajo del efecto
de estandarizar o de cambiar $k$.

\item \textbf{Mayor F1:} $k=1$, Euclidiana ($F1=0.2307$).

\item \textbf{Mayor Recall:} $k=1$, Euclidiana ($\text{Recall}=0.2144$).

\item \textbf{¿Coinciden F1 y Recall en la misma configuración?}
Sí, ambas apuntan a $k=1$ con distancia euclidiana, aunque vale la pena
matizar dicha coincidencia ya que los tres valores de $k$ comparados
aquí ya fueron preseleccionados por ranking de F1 en el Ejercicio 8, y
en este dataset el Recall resulta monótonamente decreciente en $k$ ($\uparrow \text{k}$, $\downarrow \text{Recall}$), así que el ranking de F1 y el de Recall casi
tenían que coincidir dentro de ese subconjunto de tres, por lo que no
es del todo una confirmación independiente. Dicho esto, el mecanismo
de fondo sí es real, pues a diferencia del trade-off marcado entre
$k=[1,9]$ del Ejercicio 8, aquí cambiar la métrica de distancia no
altera lo suficiente cuál vecino es ``más cercano'' como para invertir
el orden de los modelos, ya que con los datos ya estandarizados la
diferencia entre penalizar cuadráticamente (Euclidiana) o linealmente
(Manhattan) las diferencias por variable resulta menor que la que
introduce cambiar $k$.

\end{enumerate}
